\newpage

\mychapter{IV. RESULT AND DISCUSSION}{IV. RESULT AND DISCUSSION}

\section{Evaluation Results}

\hs Having described how the system was built, this chapter turns to the question that
matters most: does it actually work? To answer it, the system was tested on a manually
constructed dataset of 30 questions, spread across 10 legal categories and drawn from all
four indexed Cambodian labour laws. Three metrics were recorded for every run, namely
Retrieval Precision@5, Citation Accuracy, and Grounding Rate, each capturing a different
aspect of performance. Table~\ref{tab:overall_results} gives the headline results for the
final configuration of the system.

\begin{table}[H]
\centering
\caption{Overall Evaluation Results}
\label{tab:overall_results}
\begin{tabular}{lcc}
\toprule
\textbf{Metric} & \textbf{Score} & \textbf{Count} \\
\midrule
Retrieval Precision@5 & 76.7\% & 23/30 \\
Citation Accuracy     & 70.0\% & 21/30 \\
Grounding Rate        & 76.7\% & 23/30 \\
Average Retries per Query & 1.23 & --- \\
\bottomrule
\end{tabular}
\end{table}

\hs A single overall figure can hide a lot, so it helps to look at where the system
succeeds and where it struggles. Table~\ref{tab:category_results} breaks the results down
by category, showing both retrieval and citation performance for each of the 10 legal
topics in the dataset.

\begin{table}[H]
\centering
\caption{Per-Category Retrieval and Citation Results}
\label{tab:category_results}
\begin{tabular}{lcccc}
\toprule
\textbf{Category} & \textbf{Questions} & \textbf{Retrieval Hits} & \textbf{Precision@5} & \textbf{Citation Accuracy} \\
\midrule
Employment Contract  & 6 & 6 & 100.0\% & 100.0\% \\
Termination          & 2 & 2 & 100.0\% & 100.0\% \\
Working Hours        & 3 & 3 & 100.0\% & 100.0\% \\
Wages \& Compensation & 3 & 3 & 100.0\% & 33.3\% \\
Discrimination       & 1 & 1 & 100.0\% & 100.0\% \\
Dispute Resolution   & 1 & 1 & 100.0\% & 100.0\% \\
Trade Union          & 4 & 3 & 75.0\%  & 75.0\% \\
Leave \& Holidays    & 4 & 3 & 75.0\%  & 75.0\% \\
Workplace Regulations & 3 & 1 & 33.3\% & 33.3\% \\
Social Security      & 3 & 0 & 0.0\%  & 0.0\% \\
\midrule
\textbf{Total}       & \textbf{30} & \textbf{23} & \textbf{76.7\%} & \textbf{70.0\%} \\
\bottomrule
\end{tabular}
\end{table}

\hs These final numbers were not reached in a single attempt. The system went through
three evaluation runs, and the changes made between them had a clear, and sometimes
surprising, effect on performance. Table~\ref{tab:iteration_results} records this
progression, and the lesson behind it is discussed later in this chapter.

\begin{table}[H]
\centering
\caption{System Performance Across Evaluation Iterations}
\label{tab:iteration_results}
\begin{tabular}{lcccc}
\toprule
\textbf{Run} & \textbf{Changes Applied} & \textbf{Precision@5} & \textbf{Citation Acc.} & \textbf{Grounding Rate} \\
\midrule
Run 1 (Baseline) & Original system & 66.7\% & 66.7\% & 70.0\% \\
Run 2 & Candidate pool 20→30, aggressive rewriter & 63.3\% & 63.3\% & 73.3\% \\
Run 3 (Final) & Intent fix, grounding 200→500 chars, revised rewriter & 76.7\% & 70.0\% & 76.7\% \\
\bottomrule
\end{tabular}
\end{table}

\section{Discussion}

\subsection{Strong-Performing Categories}

\hs The clearest success is that six of the ten categories reached a perfect Precision@5
of 100\%: employment contract, termination, working hours, wages and compensation,
discrimination, and dispute resolution. Two factors explain why these topics worked so
well. The first is vocabulary. The relevant articles in the Labor Law are written in
consistent, formal Khmer that lines up well both with the terms the query rewriter
produces and with what the embedding model has seen before. The second is distinctiveness.
The core articles for these topics, such as Article 137 on working hours, Article 75 on
termination notice periods, and Article 12 on discrimination, contain content specific
enough that either FAISS or BM25 can pull them into the top-20 candidate pool, and usually
both do.

\hs The direct article lookup path deserves a separate mention, because it contributed to
perfect results on queries that name a specific article. These requests skip the retrieval
pipeline entirely and match by exact article number instead. This is a small but telling
result: it shows that intent classification earns its place in the architecture. By
recognising a lookup request and routing it to an exact-match function, the system avoids
the unnecessary imprecision of semantic search for a task that does not need it.

\subsection{Citation Gap Analysis}

\hs One result that stands out is the gap between Retrieval Precision@5, at 76.7\%, and
Citation Accuracy, at 70.0\%. The two figures are not the same thing, and the difference
between them is informative. In two queries, the system retrieved the correct article but
never referenced it with an inline \texttt{[N]} marker in the answer, so the
citation-filtering step removed it. The wages and compensation category shows this most
starkly: retrieval was perfect at 3 out of 3, yet citation accuracy was only 33.3\%. In
these cases the answer used the retrieved information correctly but described it in general
terms, without explicitly attributing it to a numbered article. The implication is
practical. The instruction in the answer prompt to mark citations needs to be firmer,
particularly for minimum-wage topics where the model tends to summarise rather than cite.

\subsection{Weak-Performing Categories}

\hs The hardest category, by a wide margin, was social security, which scored 0\% on
retrieval across all three runs. The cause is a genuine vocabulary mismatch rather than a
bug. The Social Security Law leans heavily on specialised administrative Khmer, with terms
such as {\khmerfont ហានិភ័យការងារ}, {\khmerfont ភាគទាន}, and {\khmerfont បេឡាជាតិ}, which
rarely appear in the way an ordinary person would phrase a question. Two effects compound
here. The embedding model, \texttt{gemini-embedding-001}, has limited sub-word coverage of
Khmer, so semantic matching is weak to begin with; and BM25, which relies on exact words,
fails when the user paraphrases rather than quotes the law. With neither method finding
purchase, the article is missed entirely. This is the clearest evidence in the whole
evaluation of the low-resource difficulty discussed in the literature.

\hs The workplace regulations category tells a different and more instructive story. It
fell from 3 out of 3 in the baseline to 1 out of 3 in the final run, a regression on two
questions: the disciplinary time limit in Article 26 and the union withdrawal process in
Article 7 of the Trade Union Law. Tracing the cause showed that the more aggressive query
rewriting introduced in Run 2 had injected cross-law terminology into the queries, which
pulled BM25 toward penalty articles instead of the procedural articles actually being
asked about. Run 3 corrected the rewriter and recovered most of the lost ground, but one
of the two questions never fully came back. This regression is also the reason the overall
score dipped in Run 2 before rising in Run 3, and it is a reminder that a change which
sounds like an improvement does not always behave like one.

\subsection{Grounding and Retry Analysis}

\hs The grounding rate of 76.7\% records how often the grounding node judged the retrieved
context good enough to answer from. One change made a visible difference here. Increasing
the preview given to the grounding model from 200 to 500 characters per chunk allowed it
to accept contexts it had previously rejected, since many articles open with a long
preamble that consumed most of the original 200-character window before the substantive
text began. The average of 1.23 retries per query is also worth noting: it means most
queries went through at least one rewrite-and-retry cycle before the grounding check was
satisfied. Far from being wasted effort, this confirms that the query rewriter genuinely
improves on the user's original wording, since the system frequently needed that improved
query to find relevant articles.

%\imgplaceholder[0.95\textwidth]{Evaluation Results Bar Chart — Precision@5 and Citation Accuracy per Category}{fig:eval_barchart}
\fg{0.9\textwidth}{Evaluation Results Bar Chart — Precision@5 and Citation Accuracy per Category}{images/img-ch4/Evaluation-result-chart.png}


\section{System Interface and Deployment}

\hs Beyond the numbers, the system was built into a complete, working web application and
deployed so that it could be used in practice. In production it runs as a set of Docker
containers on an Amazon EC2 instance, with Nginx serving the React frontend, forwarding
API requests to the FastAPI backend, and securing traffic over HTTPS. The figures below
show the main parts of the interface as a user encounters them.

\subsection{Welcome Screen}

\hs On first load, the application shows a welcome screen carrying the system name,
\textit{Niti}, and a four-card grid of the law categories it covers. A user can begin
either by clicking one of the suggested question cards or by typing a question straight
into the message box, so the system is approachable whether or not the user knows what to
ask.

%\imgplaceholder[0.9\textwidth]{System Welcome Screen}{fig:welcome_screen}
\fg{0.9\textwidth}{Chatbot Welcome Page}{images/img-ch4/Chatbot-Welcome-screen.png}

\subsection{Legal Question and Answer Interface}

\hs Once a question is submitted in Khmer, a typing indicator appears while the agent
pipeline does its work. When the answer is ready, it is streamed into the chat window a few
characters at a time, so the response feels live rather than appearing all at once. After
the streaming finishes, citation chips appear beneath the answer, each showing the article
number, the law name, and the page number of the source it points to.

%\imgplaceholder[0.9\textwidth]{Chat Interface with Khmer Legal Answer and Citation Chips}{fig:chat_interface}
\fg{0.9\textwidth}{Chat Interface}{images/img-ch4/Chat-Interface.png}

\subsection{PDF Viewer with Citation Navigation}

\hs Clicking a citation chip opens a PDF viewer that slides in from the right of the
screen. It renders the combined law document with pdfjs-dist and jumps straight to the
page of the cited article. Getting to the correct page required care, because the combined
PDF has a table of contents and cover pages, so the printed page numbers do not match the
digital page positions. A page-offset map handles this correction per law, so the article
the user clicked is the article they actually see.

%\imgplaceholder[0.9\textwidth]{PDF Viewer Panel Showing Cited Article}{fig:pdf_viewer}
\fg{0.9\textwidth}{Chatbot Welcome Page}{images/img-ch4/pdf-viewer.png}

\section{Summary}

\hs Taken together, the results show that the proposed Agentic RAG system works well for
its intended purpose. On a 30-question dataset spanning four Cambodian labour laws, it
reached a Retrieval Precision@5 of 76.7\% and a Citation Accuracy of 70.0\%, with the
strongest performance on the most common labour-law topics: employment contracts, working
hours, termination, and discrimination. Its clearest weakness lies in the Social Security
Law, where the gap between everyday phrasing and formal administrative Khmer defeats both
forms of retrieval, a difficulty that reflects the broader challenge of working in a
low-resource language. Alongside the measured performance, the system was realised as a
complete, deployed web application that answers questions in Khmer and links each answer
directly to its source in the original law. Both the evaluation and the working
deployment support the central claim of this thesis: that an agentic, hybrid RAG approach
is a viable way to make legal information more accessible in a low-resource setting.
